\begin{figure}[!tbp]
	\centering
	\begin{tikzpicture}[
		node distance=7mm and 8mm,
		workflowbox/.style={
			draw,
			rounded corners=2pt,
			align=center,
			minimum height=14mm,
			text width=3.7cm,
			font=\small
		},
		arrow/.style={-{Latex[length=2mm]}, thick}
	]
		\node[workflowbox] (generation) {
			\textbf{1. Generate agglomerates}\\
			size, compactness and\\random realizations\\DEM consistency check
		};
		\node[workflowbox, right=of generation] (targets) {
			\textbf{2. Define target}\\
			mean geometric\\exposure
		};
		\node[workflowbox, right=of targets] (graphs) {
			\textbf{3. Build graphs}\\
			particles as nodes and\\contacts as edges
		};
		\node[workflowbox, below=of graphs] (compression) {
			\textbf{4. Learn descriptor}\\
			graph $\rightarrow \boldsymbol{\lambda}_A$\\$\rightarrow$ exposure
		};
		\node[workflowbox, left=of compression] (interpretation) {
			\textbf{5. Interpret and compare}\\
			physical meaning and\\comparison with scalars
		};

		\draw[arrow] (generation) -- (targets);
		\draw[arrow] (targets) -- (graphs);
		\draw[arrow] (graphs) -- (compression);
		\draw[arrow] (compression) -- (interpretation);
	\end{tikzpicture}
	\caption{Five-step workflow for generating synthetic agglomerates, representing them as graphs, learning a compact structural descriptor and evaluating its physical meaning and predictive value. DEM is an auxiliary consistency check within generation, not a separate characterization step. The learning and comparison stages remain pending for the new dataset.}
	\label{fig:workflow}
\end{figure}
